\documentclass{article}
\usepackage[utf8]{inputenc}
\usepackage{hyperref}
\usepackage{geometry}
\geometry{a4paper, margin=1in}

\title{Scientific Papers \& Observational Backing \\ \large Superluminal Astrophysical Phenomena (2024-2025 Research)}
\author{Solar System Simulation Research}
\date{February 24, 2026}

\begin{document}

\maketitle

This document archives the cutting-edge scientific papers, reviews, and observational data used to form the computational and mathematical basis of our SMBH shockwave models.

\section*{Area 1: General Relativistic Magnetohydrodynamics (GRMHD)}

\textbf{1. "Hybrid Black Hole-Disk Driven Jets: Steady Axisymmetric Ideal MHD Modeling" (2025)}
\begin{itemize}
    \item \textbf{Summary:} Describes tractable theoretical models for simulating relativistic jets driven by both the black hole's spin and the accretion disk.
    \item \textbf{Observational Link:} Explains the energy distribution and Lorentz factors observed in jets like M87 and Cen A.
\end{itemize}

\textbf{2. "Resistive Relativistic MHD Simulations of Astrophysical Jets" (2023-2024)}
\begin{itemize}
    \item \textbf{Summary:} Investigates the impact of plasma resistivity on turbulence, current sheet formation, and magnetic reconnection plasmoids within jets.
    \item \textbf{Observational Link:} Helps explain the rapid, localized "flaring" events seen in Sgr A*.
\end{itemize}

\textbf{3. "Superluminal Motion in Jet of NLSy1 Galaxy TXS 1206+549" (2024)}
\begin{itemize}
    \item \textbf{Summary:} Uses VLBI data to track and calculate superluminal motion parameters in a narrow-line Seyfert 1 galaxy jet.
    \item \textbf{Observational Link:} Direct mathematical application of the apparent superluminal velocity equation $v_{app} = \frac{v \sin\theta}{1 - (v/c)\cos\theta}$.
\end{itemize}

\section*{Area 2: Particle-In-Cell (PIC) Plasma Simulations}

\textbf{4. "Superluminal Wave Activation at Collisionless Shocks" (2024)}
\begin{itemize}
    \item \textbf{Summary:} PIC simulations validating mechanisms where Alfvénic perturbations at relativistic shocks convert into superluminal O-modes.
    \item \textbf{Observational Link:} Proposed mechanism for Fast Radio Bursts (FRBs).
\end{itemize}

\textbf{5. "Kinetic Instabilities in Relativistic Jets: Weibel and Mushroom"}
\begin{itemize}
    \item \textbf{Summary:} Uses PIC codes to model the microscopic interactions of electron-positron/electron-ion plasmas.
    \item \textbf{Observational Link:} Explains the origin of the $10^{15}$ Gauss magnetic fields required for the Magnetar Relay (H11).
\end{itemize}

\section*{Area 3: Numerical Relativity \& Gravitational Wave Memory}

\textbf{6. "The Cauchy-Characteristic Evolution (CCE) Method for GW Memory" (2024)}
\begin{itemize}
    \item \textbf{Summary:} Details the mathematical CCE method for extracting highly accurate gravitational waveforms at null infinity, fully capturing the Christodoulou memory effect.
    \item \textbf{Observational Link:} Essential for tracking the "Cosmic Ray Surfboard" (H8) hypothesis.
\end{itemize}

\textbf{7. "Machine Learning Surrogate Models for GW Emulating" (2025)}
\begin{itemize}
    \item \textbf{Summary:} The SXS group's development of ML-driven surrogate models to rapidly interpolate between expensive numerical relativity simulations.
    \item \textbf{Observational Link:} Bridges the gap between theoretical Einstein equations and massive data streams from next-generation detectors.
\end{itemize}

\end{document}
